\documentclass{article}

\usepackage{amsmath,amssymb}
\usepackage{xurl}
\usepackage[hidelinks]{hyperref}
\setlength{\emergencystretch}{2em}

% Shared commands for notes in docs/paper-gaps/.

\DeclareUnicodeCharacter{2080}{\ifmmode{}_{0}\else\textsubscript{0}\fi}
\DeclareUnicodeCharacter{2081}{\ifmmode{}_{1}\else\textsubscript{1}\fi}
\DeclareUnicodeCharacter{2082}{\ifmmode{}_{2}\else\textsubscript{2}\fi}
\DeclareUnicodeCharacter{2099}{\ifmmode{}_{n}\else\textsubscript{n}\fi}

\newcommand{\Lean}{\textsc{Lean}}
\ExplSyntaxOn
\NewDocumentCommand{\leanid}{m}{
  \ifmmode
    \textsf{\detokenize{#1}}
  \else
    \group_begin:
    \urlstyle{sf}
    \tl_set:Nn \l_tmpa_tl {#1}
    \tl_replace_all:Nnn \l_tmpa_tl {\_} {_}
    \exp_args:NV \path \l_tmpa_tl
    \group_end:
  \fi
}
\ExplSyntaxOff
\providecommand{\tr}{\operatorname{tr}}
\newcommand{\tnleanIssueBase}{https://github.com/LionSR/TNLean/issues/}
\newcommand{\tnleanPrBase}{https://github.com/LionSR/TNLean/pull/}
\newcommand{\tnleanDocBase}{https://sirui-lu.com/TNLean/docs/}
\newcommand{\ghissue}[1]{\href{\tnleanIssueBase#1}{GitHub issue \##1}}
\newcommand{\ghpr}[1]{\href{\tnleanPrBase#1}{PR \##1}}
\newcommand{\doclink}[1]{\href{\tnleanDocBase#1.html}{\path{#1}}}

% Dirac notation and the identity operator, as in blueprint/src/macros/common.tex.
% \providecommand keeps note-local definitions authoritative.
\usepackage{dsfont}
\providecommand{\ket}[1]{|#1\rangle}
\providecommand{\bra}[1]{\langle #1|}
\providecommand{\braket}[2]{\langle #1 | #2 \rangle}
\providecommand{\ketbra}[2]{|#1\rangle\!\langle #2|}
\providecommand{\Id}{\mathds{1}}
\providecommand{\one}{\mathds{1}}
\providecommand{\C}{\mathbb{C}}
\providecommand{\MN}[1]{M_{#1}(\C)}
\providecommand{\MD}{M_D(\C)}

% External referencing for paper sources.  The build helper writes sanitized
% auxiliary files under build/paper-aux/ and excludes duplicated source labels.
\usepackage{xr}

% Import a generated paper auxiliary file with a caller-supplied label prefix.
% The candidate paths support builds from docs/paper-gaps/, the repository root,
% and build/paper-aux/ itself.
\newcommand{\importpaperaux}[2]{%
  \IfFileExists{../../build/paper-aux/#2.aux}{%
    \externaldocument[#1]{../../build/paper-aux/#2}%
  }{%
    \IfFileExists{build/paper-aux/#2.aux}{%
      \externaldocument[#1]{build/paper-aux/#2}%
    }{%
      \IfFileExists{#2.aux}{%
        \externaldocument[#1]{#2}%
      }{}%
    }%
  }%
}

% General external reference macros
\newcommand{\paperref}[2]{\ref{#1-#2}}
\newcommand{\papereqref}[2]{(\ref{#1-#2})}

% CPSV16 source (1606.00608 / MPDO-22-12-17-2.tex) external document and shortcuts
\importpaperaux{cpsv16-}{1606.00608/MPDO-22-12-17-2-tnlean}
\newcommand{\cpsvref}[1]{\paperref{cpsv16}{#1}}
\newcommand{\cpsveqref}[1]{\papereqref{cpsv16}{#1}}

% Verdict marker: \gapnote{<kind>}{<status>}. Kind is the policy
% classification (clarification, local-correction, scope-restriction,
% unfaithful, false-source, open-gap); status is open, wip (elimination
% underway), resolved, or historical. Machine-read by the site index;
% prints a verdict line.
\newcommand{\gapnote}[2]{\par\noindent\textbf{Verdict:} #1 (#2).\par}


\title{Scalar Normalization in Wolf's Lorentz Normal Form}
\author{}
\date{2026-07-16}

\begin{document}
\maketitle

\section{The assertion}

Wolf's Proposition~2.11 states that every qubit channel
$T:\MN{2}\to\MN{2}$ admits invertible completely positive maps
$\Phi_1,\Phi_2$, each of Kraus rank one, such that
\begin{align}
  T' = \Phi_2\circ T\circ\Phi_1
\end{align}
is again a qubit channel and belongs to one of three Lorentz normal forms:
diagonal, non-diagonal, or singular.\footnote{\cite[Proposition~2.11]{Wolf2012Quantum};
local source \path{Notes/WolfNoteTexSource/ch02_representations.tex},
lines 1021--1035.}  Each filtering has the form
\begin{align}
  \Phi_X(A)=XAX^\dagger,
  \qquad X\in\mathrm{GL}(2,\C).
\end{align}
Thus the proposition allows every invertible Kraus-rank-one completely positive
filter, not only those represented by determinant-one matrices.

\section{Determinant-one action and scalar freedom}

Choose a square root $c^2=\det X$ and write $X=cS$, where
$S\in\mathrm{SL}(2,\C)$. Then
\begin{align}
  \Phi_X(A)
    &= |c|^2\Phi_S(A),\\
  |c|^2
    &= |\det X|.
\end{align}
Consequently, if $X_i=c_iS_i$ for the two filters, then
\begin{align}
  \Phi_{X_2}\circ T\circ\Phi_{X_1}
  = |c_1c_2|^2
    (\Phi_{S_2}\circ T\circ\Phi_{S_1}).
  \label{eq:scalar-factor}
\end{align}
The matrices $S_1,S_2$ determine the determinant-one Lorentz action. The
positive factor $|c_1c_2|^2$ selects the normalization on the resulting ray of
completely positive maps.

This distinction is immaterial when only the Lorentz orbit up to positive scale
is considered. It is essential when the representative is required to be a
channel. If the determinant-one representative satisfies
\begin{align}
  \tr\!\left(\Phi_{S_2}\circ T\circ\Phi_{S_1}(A)\right)
  = \alpha\,\tr(A),
\end{align}
then it is trace preserving only when $\alpha=1$. General invertible filters
can multiply the map by $\alpha^{-1}$ through the scalar in
\eqref{eq:scalar-factor}; determinant-one filters have no such freedom. The
minimisation in Propositions~2.8 and~2.9 produces partial traces proportional to
the identity, but the proportionality scalar must still be chosen so that the
qubit representative is trace preserving.

\section{The point at issue}

A former local statement required both filters to be determinant-one while also
requiring the filtered map to be a qubit channel in one of the three canonical
forms.\footnote{The removed declaration was
\leanid{Wolf.exists_lorentz_normal_form_qubit} in
\doclink{TNLean/Channel/LorentzNormalForm}.}  Equation~\eqref{eq:scalar-factor}
shows why the determinant-one restriction does not justify the normalization in
Wolf's proposition. The source theorem permits general invertible filters, and
its channel conclusion uses their scalar freedom. The restricted local
statement was therefore removed rather than treated as a formalization of
Proposition~2.11.

The three canonical-form predicates and the square full-Kraus-rank minimisation
theorem are formalized.\footnote{Formal declarations
\leanid{Wolf.IsLorentzDiagonal}, \leanid{Wolf.IsLorentzNonDiagonal},
\leanid{Wolf.IsLorentzSingular}, and
\leanid{Wolf.exists_normal_form_generic} in
\doclink{TNLean/Channel/LorentzNormalForm}.}  The theorem asserting existence of
a canonical representative for every qubit channel remains open.

\section{Elimination plan}

A source-faithful proof requires the following steps.
\begin{enumerate}
  \item Define general invertible Kraus-rank-one completely positive filters and
    separate each filter into a determinant-one factor and a nonzero scalar.
  \item Use the minimisation argument of Propositions~2.8 and~2.9 for the
    full-Kraus-rank case, retaining the proportionality constants of the two
    marginals.
  \item Formalize the Lorentz action of $\mathrm{SL}(2,\C)$ on the Pauli
    transfer matrix and classify its diagonal, non-diagonal, and singular orbits.
  \item Choose the scalar factors so that the canonical representative is trace
    preserving, and verify complete positivity and the stated parameter bounds.
  \item Prove the three-case existence theorem using the general filters and link
    it to the existing canonical-form predicates.
\end{enumerate}

The discrepancy is an open gap: the canonical forms are defined, and the generic
minimisation theorem is available in square dimension, but the scalar-normalized
Lorentz-orbit classification required by Proposition~2.11 has not yet been
formalized.

\bibliographystyle{alpha}
\bibliography{references}

\end{document}